\documentclass[a4paper,landscape,8pt]{extarticle}
\usepackage[landscape,margin=0.7cm,top=0.5cm,bottom=0.7cm]{geometry}
\usepackage[utf8]{inputenc}
\usepackage[T1]{fontenc}
\usepackage[ngerman]{babel}
\usepackage{helvet}
\renewcommand{\familydefault}{\sfdefault}
\usepackage{xcolor}
\usepackage{tikz}
\usetikzlibrary{positioning}
\usepackage{enumitem}
\usepackage{multicol}
\usepackage{array}
\usepackage{fancyhdr}
\usepackage{amssymb}

% Farben
\definecolor{wagreen}{HTML}{25D366}
\definecolor{chatblue}{HTML}{1565C0}
\definecolor{warnred}{HTML}{C62828}
\definecolor{lightgray}{HTML}{F5F5F5}
\definecolor{bordergray}{HTML}{BBBBBB}

\pagestyle{fancy}
\fancyhf{}
\renewcommand{\headrulewidth}{0pt}
\fancyfoot[C]{\footnotesize Seite \thepage{} von 3}

\setlength{\parindent}{0pt}
\setlength{\parskip}{2pt}

% Kompakte Listen
\setlist{nosep, leftmargin=1em, itemsep=0pt, topsep=1pt}

\begin{document}

% ============================================================================
% SEITE 1: INFORMATIONSBLATT
% ============================================================================

\noindent{\Large\bfseries WhatsApp \& Co.: Schnell kommunizieren} \hfill
{\footnotesize\textbf{Lernziele:} Nachrichten senden | Fotos teilen | Gefahren erkennen}

\vspace{1mm}
\hrule
\vspace{1mm}

% Drei-Spalten-Layout
\begin{minipage}[t]{0.32\linewidth}
\begin{center}
\colorbox{wagreen!15}{%
\begin{minipage}{0.95\linewidth}
\vspace{2mm}
\begin{center}
{\Large\bfseries\textcolor{wagreen}{1. Was ist WhatsApp?}}\\[1pt]
{\footnotesize\itshape Mehr als SMS}
\end{center}
\vspace{2mm}
\end{minipage}%
}
\end{center}

\vspace{1mm}

\textbf{Ein Messenger ist:}
\begin{itemize}
\item Eine App für \textbf{schnelle Nachrichten}
\item Funktioniert über \textbf{Internet} (WLAN/mobile Daten)
\item \textbf{Kostenlos} -- keine SMS-Gebühren
\end{itemize}

\textbf{WhatsApp vs. SMS:}
\begin{itemize}
\item SMS: Pro Nachricht zahlen, nur Text
\item WhatsApp: Kostenlos, Fotos, Videos, Sprachnachrichten
\end{itemize}

\textbf{Analogie:} Wie ein \textbf{Gespräch} -- schnell hin und her, statt formeller Briefe.

\textbf{Die drei Bereiche:}
\begin{itemize}
\item \textbf{Chats:} Deine Gespräche
\item \textbf{Status:} Bilder, die nach 24h verschwinden
\item \textbf{Anrufe:} Telefonate über Internet
\end{itemize}

\end{minipage}%
\hfill%
\begin{minipage}[t]{0.32\linewidth}
\begin{center}
\colorbox{chatblue!10}{%
\begin{minipage}{0.95\linewidth}
\vspace{2mm}
\begin{center}
{\Large\bfseries\textcolor{chatblue}{2. Nachrichten \& Fotos}}\\[1pt]
{\footnotesize\itshape So geht's}
\end{center}
\vspace{2mm}
\end{minipage}%
}
\end{center}

\vspace{1mm}

\textbf{Nachricht senden:}
\begin{enumerate}
\item WhatsApp öffnen
\item Chat mit Person antippen
\item Unten ins \textbf{Textfeld} tippen
\item Text eingeben
\item \textbf{Grüner Pfeil} = Senden
\end{enumerate}

\textbf{Foto senden:}
\begin{enumerate}
\item Im Chat: \textbf{+} links antippen
\item ``Foto- und Videomediathek''
\item Foto auswählen
\item Optional: Text dazu schreiben
\item Senden
\end{enumerate}

\textbf{Die Haken verstehen:}
\begin{itemize}
\item $\checkmark$ \textbf{Ein} grauer Haken: gesendet
\item $\checkmark\checkmark$ \textbf{Zwei} graue Haken: angekommen
\item \textcolor{chatblue}{$\checkmark\checkmark$} \textbf{Zwei blaue} Haken: gelesen
\end{itemize}

\end{minipage}%
\hfill%
\begin{minipage}[t]{0.32\linewidth}
\begin{center}
\colorbox{warnred!10}{%
\begin{minipage}{0.95\linewidth}
\vspace{2mm}
\begin{center}
{\Large\bfseries\textcolor{warnred}{3. Vorsicht: Betrug!}}\\[1pt]
{\footnotesize\itshape Der ``Enkel-Trick'' digital}
\end{center}
\vspace{2mm}
\end{minipage}%
}
\end{center}

\vspace{1mm}

\fcolorbox{warnred}{warnred!5}{%
\begin{minipage}{0.92\linewidth}
\footnotesize
\textbf{Typische Betrugsmasche:}\\
``Hallo Oma, ich habe eine neue Nummer. Mein Handy ist kaputt. Kannst du mir schnell Geld überweisen?''
\end{minipage}%
}

\vspace{2mm}

\textbf{So schützen Sie sich:}
\begin{itemize}
\item \textbf{IMMER} auf alter Nummer anrufen und nachfragen!
\item Keine Geld-Überweisungen an ``neue'' Nummern
\item Im Zweifel: Auflegen, selbst anrufen
\end{itemize}

\textbf{Verdächtige Zeichen:}
\begin{itemize}
\item Unbekannte Nummer
\item Dringend Geld gebraucht
\item Bitte um Geheimhaltung
\item Keine Rückruf-Möglichkeit
\end{itemize}

\textbf{Regel:} Wer wirklich Familie ist, kann auch auf der \textbf{alten Nummer} erreicht werden!

\end{minipage}

\vspace{1mm}

% Zusammenfassung
\begin{center}
\fcolorbox{bordergray}{lightgray}{%
\begin{minipage}{0.98\linewidth}
\centering
\vspace{1.5mm}
\textbf{Zusammenfassung:}
\textcolor{wagreen}{\textbf{WhatsApp}} = kostenlos über Internet ~$|$~
\textcolor{chatblue}{\textbf{Blaue Haken}} = gelesen ~$|$~
\textcolor{warnred}{\textbf{``Neue Nummer''}} = erst anrufen und prüfen!
\vspace{1.5mm}
\end{minipage}%
}
\end{center}

\newpage

% ============================================================================
% SEITE 2: ÜBUNGEN
% ============================================================================

\begin{center}
{\LARGE\bfseries Übungen: WhatsApp sicher nutzen}\\[3pt]
{\normalsize Schau dir Seite 1 an und beantworte die Fragen}
\end{center}

\vspace{3mm}
\hrule
\vspace{4mm}

\begin{multicols}{2}

\textbf{\large Teil A: Grundlagen verstehen}

Richtig oder Falsch? Kreuze an.

\vspace{2mm}

\renewcommand{\arraystretch}{1.6}
\begin{tabular}{@{}p{8cm}cc@{}}
& R & F \\
1. WhatsApp-Nachrichten kosten wie SMS. & $\square$ & $\square$ \\
2. WhatsApp braucht Internet (WLAN oder mobile Daten). & $\square$ & $\square$ \\
3. Ein blauer Haken bedeutet: Nachricht gesendet. & $\square$ & $\square$ \\
4. Zwei blaue Haken bedeuten: Nachricht gelesen. & $\square$ & $\square$ \\
5. Status-Bilder bleiben für immer sichtbar. & $\square$ & $\square$ \\
\end{tabular}
\renewcommand{\arraystretch}{1}

\vspace{6mm}

\textbf{\large Teil B: Die Haken}

Was bedeuten die Haken? Verbinde.

\vspace{2mm}

\renewcommand{\arraystretch}{1.8}
\begin{tabular}{@{}p{3cm}p{4cm}@{}}
$\checkmark$ & $\circ$ ~~~~ $\circ$ Nachricht wurde gelesen \\
$\checkmark\checkmark$ (grau) & $\circ$ ~~~~ $\circ$ Nachricht wurde gesendet \\
$\checkmark\checkmark$ (blau) & $\circ$ ~~~~ $\circ$ Nachricht ist angekommen \\
\end{tabular}
\renewcommand{\arraystretch}{1}

\vspace{6mm}

\textbf{\large Teil C: Foto senden}

Nummeriere die Schritte in der richtigen Reihenfolge (1--5):

\vspace{2mm}

\renewcommand{\arraystretch}{1.8}
\begin{tabular}{@{}cp{7cm}@{}}
$\square$ & Foto aus der Galerie auswählen \\
$\square$ & Chat mit der Person öffnen \\
$\square$ & Auf den Senden-Pfeil tippen \\
$\square$ & Das + Symbol antippen \\
$\square$ & WhatsApp öffnen \\
\end{tabular}
\renewcommand{\arraystretch}{1}

\columnbreak

\textbf{\large Teil D: Betrug erkennen}

Lesen Sie die Nachrichten. Betrug oder echt? Was tun Sie?

\vspace{3mm}

\textbf{Nachricht 1:}\\
``Hallo Oma, mein Handy ist kaputt. Das ist meine neue Nummer. Kannst du mir 500 Euro überweisen? Dringend!''

\vspace{2mm}
$\square$ Echt ~~~~~ $\square$ Betrug

\textit{Was tun Sie?}
\vspace{4mm}

\rule{\linewidth}{0.4pt}

\vspace{5mm}

\textbf{Nachricht 2:}\\
Von Ihrer Tochter (bekannte Nummer): ``Mama, ich schicke dir gleich Fotos vom Geburtstag!''

\vspace{2mm}
$\square$ Echt ~~~~~ $\square$ Betrug

\textit{Warum?}
\vspace{4mm}

\rule{\linewidth}{0.4pt}

\vspace{5mm}

\textbf{Nachricht 3:}\\
``Sie haben gewonnen! Klicken Sie hier für Ihren Preis: www.gewinn-jetzt.xyz''

\vspace{2mm}
$\square$ Echt ~~~~~ $\square$ Betrug

\textit{Was tun Sie?}
\vspace{4mm}

\rule{\linewidth}{0.4pt}

\end{multicols}

\newpage

% ============================================================================
% SEITE 3: CHECKLISTE
% ============================================================================

\begin{center}
{\LARGE\bfseries Checkliste: Das habe ich verstanden}\\[3pt]
{\normalsize Geh die Liste durch und hake ab, was du sicher weißt}
\end{center}

\vspace{3mm}
\hrule
\vspace{6mm}

% Drei Boxen nebeneinander
\begin{minipage}[t]{0.31\linewidth}
\begin{center}
\fcolorbox{wagreen}{wagreen!15}{%
\begin{minipage}{0.92\linewidth}
\vspace{3mm}
\begin{center}
{\large\bfseries\textcolor{wagreen}{WhatsApp-Grundlagen}}
\end{center}
\vspace{2mm}
\begin{itemize}[label=$\square$, itemsep=5pt]
\item Ich weiß, dass WhatsApp \textbf{Internet} braucht
\item Ich kann die \textbf{App öffnen}
\item Ich finde meine \textbf{Chats}
\item Ich verstehe die \textbf{drei Bereiche}
\end{itemize}
\vspace{3mm}
\end{minipage}%
}
\end{center}
\end{minipage}%
\hfill%
\begin{minipage}[t]{0.31\linewidth}
\begin{center}
\fcolorbox{chatblue}{chatblue!8}{%
\begin{minipage}{0.92\linewidth}
\vspace{3mm}
\begin{center}
{\large\bfseries\textcolor{chatblue}{Nachrichten \& Fotos}}
\end{center}
\vspace{2mm}
\begin{itemize}[label=$\square$, itemsep=5pt]
\item Ich kann eine \textbf{Nachricht senden}
\item Ich kann ein \textbf{Foto senden}
\item Ich verstehe die \textbf{Haken-Bedeutung}
\item Ich weiß, wie ich \textbf{antworte}
\end{itemize}
\vspace{3mm}
\end{minipage}%
}
\end{center}
\end{minipage}%
\hfill%
\begin{minipage}[t]{0.31\linewidth}
\begin{center}
\fcolorbox{warnred}{warnred!8}{%
\begin{minipage}{0.92\linewidth}
\vspace{3mm}
\begin{center}
{\large\bfseries\textcolor{warnred}{Sicherheit}}
\end{center}
\vspace{2mm}
\begin{itemize}[label=$\square$, itemsep=5pt]
\item Ich kenne den \textbf{Enkel-Trick}
\item Ich \textbf{rufe zurück} bei ``neuer Nummer''
\item Ich klicke \textbf{nicht} auf unbekannte Links
\item Ich überweise \textbf{kein Geld} ohne Prüfung
\end{itemize}
\vspace{3mm}
\end{minipage}%
}
\end{center}
\end{minipage}

\vspace{8mm}

% Gesamtverständnis
\begin{center}
\fcolorbox{black}{lightgray}{%
\begin{minipage}{0.95\linewidth}
\vspace{4mm}
\begin{center}
{\large\bfseries Gesamtverständnis}
\end{center}
\vspace{3mm}
\begin{multicols}{2}
\begin{itemize}[label=$\square$, itemsep=6pt]
\item Ich kann \textbf{Nachrichten} schreiben und senden
\item Ich kann \textbf{Fotos} aus meiner Galerie teilen
\item Ich verstehe, was die \textbf{Haken} bedeuten
\item Ich erkenne \textbf{Betrugsversuche}
\item Ich weiß, wie ich mich \textbf{schützen} kann
\item Ich fühle mich \textbf{sicher} bei WhatsApp
\end{itemize}
\end{multicols}
\vspace{3mm}
\end{minipage}%
}
\end{center}

\vspace{8mm}

% Notfall-Notizen
\begin{center}
\fcolorbox{bordergray}{white}{%
\begin{minipage}{0.95\linewidth}
\vspace{4mm}
\begin{center}
{\large\bfseries Meine Sicherheits-Notizen}\\[3pt]
{\footnotesize (Fülle das hier aus und bewahre es auf!)}
\end{center}
\vspace{4mm}
\renewcommand{\arraystretch}{2}
\begin{tabular}{@{}p{0.45\linewidth}p{0.45\linewidth}@{}}
\textbf{Telefonnummern meiner Familie:} & \textbf{Bei Verdacht wende ich mich an:} \\
\dotfill & \dotfill \\[3mm]
\dotfill & \\[3mm]
\dotfill & \\
\end{tabular}
\renewcommand{\arraystretch}{1}
\vspace{4mm}
\end{minipage}%
}
\end{center}

\vfill

\begin{center}
\footnotesize\textit{Stand: \today{} -- Bei ``Neue Nummer''-Nachrichten: IMMER erst auf der alten Nummer anrufen!}
\end{center}

\end{document}
